\notechapter{N-20220320}{Mean Value Theorems: How Local Slopes Control an Interval}

A derivative is information at one point, whereas an estimate such as
\[
  |f(b)-f(a)|\le M|b-a|
\]
compares two distinct points. The mean value theorem explains how to pass between these scales. We follow the classical route through extrema, Rolle's theorem, and its two-function extension, as developed in Rudin's \emph{Principles of Mathematical Analysis}, Chapter 5, and the MIT 18.100B real-analysis lectures.\footnote{\url{https://ocw.mit.edu/courses/18-100b-real-analysis-spring-2025/pages/lecture-notes/}}

\section{Why a compact interval enters the proof}

Suppose \(h:[a,b]\to\mathbb R\) is continuous and \(h(a)=h(b)\). The extreme value theorem guarantees that \(h\) attains a maximum and a minimum. If both extrema occur only at endpoints, their common endpoint value makes \(h\) constant. Otherwise one extremum occurs at an interior point \(c\). Provided \(h\) is differentiable there, the two one-sided difference quotients have opposite allowable signs, so \(h'(c)=0\).

\begin{theorem}[Rolle]
If \(h\) is continuous on \([a,b]\), differentiable on \((a,b)\), and \(h(a)=h(b)\), then \(h'(c)=0\) for some \(c\in(a,b)\).
\end{theorem}

The hypotheses cannot be omitted casually. For \(h(x)=|x|\) on \([-1,1]\), the endpoints agree but the sole interior minimum occurs at a point where the derivative does not exist. If continuity fails at an endpoint, the extreme value argument also breaks.

\section{Removing the secant line}

\begin{theorem}[Mean value theorem]
If \(f\) is continuous on \([a,b]\) and differentiable on \((a,b)\), then there is \(c\in(a,b)\) with
\[
  f(b)-f(a)=f'(c)(b-a).
\]
\end{theorem}
\begin{proof}
The affine interpolant between the endpoint values is
\[
  \ell(x)=f(a)+\frac{f(b)-f(a)}{b-a}(x-a).
\]
The difference \(h=f-\ell\) satisfies \(h(a)=h(b)=0\). Rolle's theorem supplies \(c\) with
\[
  0=h'(c)=f'(c)-\frac{f(b)-f(a)}{b-a}.
\]
\end{proof}

This proof suggests the correct geometric operation: subtract the simplest function matching the boundary data. For example, \(f(x)=x^2\) on \([1,3]\) has secant slope \(4\); the equality \(2c=4\) identifies \(c=2\).

\section{The two-function theorem and a nonzero denominator}

\begin{theorem}[Cauchy's mean value theorem]
Suppose \(f,g\) are continuous on \([a,b]\) and differentiable on \((a,b)\). Then some \(c\in(a,b)\) satisfies
\[
  (f(b)-f(a))g'(c)=(g(b)-g(a))f'(c).
\]
\end{theorem}
\begin{proof}
Put \(A=f(b)-f(a)\), \(B=g(b)-g(a)\), and set
\[
  h(x)=Bf(x)-Ag(x).
\]
Its endpoint difference is \(BA-AB=0\). Applying Rolle to \(h\) yields \(Bf'(c)-Ag'(c)=0\).
\end{proof}

The cross-multiplied identity is the theorem's basic form. To write a ratio one must justify the divisions. If \(g'\ne0\) throughout \((a,b)\), Rolle applied to \(g\) forces \(g(b)\ne g(a)\), and Cauchy's identity becomes
\[
  \frac{f(b)-f(a)}{g(b)-g(a)}
  =\frac{f'(c)}{g'(c)}.
\]
Without this extra condition, a quotient formulation may be meaningless although the cross-multiplied theorem remains true.

\section{Where l'Hopital's rule comes from}

Consider the right-hand \(0/0\) situation. Assume \(f,g\) are continuous on \([a,a+\delta)\), differentiable on \((a,a+\delta)\), with
\[
  f(a)=g(a)=0,\qquad g'(x)\ne0\quad(a<x<a+\delta).
\]
For every \(x>a\), Cauchy's theorem on \([a,x]\) supplies \(c_x\in(a,x)\) such that
\[
  \frac{f(x)}{g(x)}=\frac{f'(c_x)}{g'(c_x)}.
\]
Consequently, if \(f'(t)/g'(t)\to L\) as \(t\downarrow a\), then \(f(x)/g(x)\to L\), because \(c_x\downarrow a\) in the sense that \(a<c_x<x\to a\).

\begin{example}
Let \(f(x)=\sin x-x\) and \(g(x)=x^3\), with \(a=0\). Direct substitution gives \(0/0\), but
\[
 \frac{f'(x)}{g'(x)}
 =\frac{\cos x-1}{3x^2}\longrightarrow-\frac16.
\]
The derivative ratio here still uses the familiar limit \((1-\cos x)/x^2\to\tfrac12\); it does not license a circular application of l'Hopital to itself.
\end{example}

The theorem is a statement about limits under explicit hypotheses. Taking derivatives of numerator and denominator is not a general algebraic operation on fractions.

\section{Derivative bounds give global estimates}

\begin{proposition}[Lipschitz estimate]
If \(f\) is continuous on an interval and differentiable in its interior, with \(|f'(x)|\le M\), then
\[
  |f(x)-f(y)|\le M|x-y|
\]
for all \(x,y\) in the interval.
\end{proposition}
\begin{proof}
Apply the mean value theorem to the subinterval with endpoints \(x,y\), then take absolute values.
\end{proof}

If \(f'=0\) throughout an interval, this forces \(f\) to be constant. If \(f'\ge0\), the same reasoning shows \(f\) is nondecreasing. These are stronger conclusions than a picture of tangent lines: a pointwise differential inequality implies a global comparison.

The Lipschitz condition also enters uniqueness arguments for ordinary differential equations. When \(|F(t,u)-F(t,v)|\le L|u-v|\), two solutions with the same initial value satisfy
\[
  e(t):=|u(t)-v(t)|\le L\int_{t_0}^{t}e(s)\,ds
  \qquad (t\ge t_0).
\]
For \(H(t)=\int_{t_0}^t e(s)\,ds\), one has \(H(t_0)=0\) and \(H'(t)\le LH(t)\). Therefore
\((e^{-Lt}H(t))'\le0\). Since \(H\ge0\), it follows that \(H=0\) and the solutions coincide. This is the elementary core of the Gronwall argument.

\section{Continuity versus uniform continuity}

Continuity at a point \(x\) allows the \(\delta\) in the \(\varepsilon\)--\(\delta\) definition to depend on \(x\). Uniform continuity asks for the same \(\delta(\varepsilon)\) across the domain.

\begin{theorem}[Heine--Cantor]
A continuous function from a compact metric space to a metric space is uniformly continuous.
\end{theorem}
\begin{proof}
Suppose otherwise. There is \(\varepsilon_0>0\) and sequences \(x_n,y_n\) with
\[
  d(x_n,y_n)<1/n,\qquad
  d(f(x_n),f(y_n))\ge\varepsilon_0.
\]
By compactness, a subsequence \(x_{n_j}\) converges to some \(x\). The triangle inequality gives \(y_{n_j}\to x\) as well. Continuity at \(x\) forces both image subsequences to converge to \(f(x)\), a contradiction.
\end{proof}

Compactness matters. The function \(x\mapsto x^2\) on \(\mathbb R\) is continuous but not uniformly continuous: take \(x_n=n\), \(y_n=n+1/n\). The inputs approach one another, while the output difference tends to \(2\).

\section{Taylor's formula as an error statement}

For \(f\in C^2\) between \(a\) and \(x\), Taylor's theorem gives some \(\xi\) between them such that
\[
  f(x)=f(a)+f'(a)(x-a)+\frac{f''(\xi)}2(x-a)^2.
\]
In particular, if \(|f''|\le K\), the error of the tangent-line approximation satisfies
\[
  |f(x)-f(a)-f'(a)(x-a)|\le \frac K2|x-a|^2.
\]
The logical progression is now visible: extrema force a zero derivative; removing a secant produces a mean slope; bounding derivatives controls finite differences; controlling a higher derivative quantifies the error of a local polynomial approximation.
